\documentclass[11pt, a4paper]{article}

\usepackage[utf8]{inputenc}
\usepackage[T1]{fontenc}
\usepackage{amsmath, amssymb, amsfonts}
\usepackage{graphicx}
\usepackage[margin=1in]{geometry}
\usepackage{hyperref}
\usepackage[numbers, sort&compress]{natbib}
\usepackage{booktabs}
\usepackage{siunitx}
\usepackage{caption}
\usepackage{float}
\usepackage{xcolor}

\title{\textbf{Hybrid Universal Control in a 2x2 Qubit+Cavity Logical Subspace: Selective vs Native Gate Families in \texttt{cqed\_sim}}}
\author{Codex \\ cQED Autonomous Research Platform}
\date{\today}

\begin{document}

\maketitle

\begin{abstract}
We study a practical control-design question for circuit quantum electrodynamics:
which gate family is the best realistic route to hybrid universal control of a
qubit and cavity when the logical task is restricted to the smallest nontrivial
\(2\times 2\) block? The benchmark is implemented on top of \texttt{cqed\_sim}
for a dispersive transmon--storage system with
\(\omega_q/2\pi=\SI{6.150}{GHz}\), \(\omega_c/2\pi=\SI{5.241}{GHz}\),
\(\alpha/2\pi=\SI{-255}{MHz}\), \(\chi/2\pi=\SI{-2.84}{MHz}\),
\(\chi'/2\pi=\SI{-21}{kHz}\), and \(K/2\pi=\SI{-28}{kHz}\). We compare six
routes: (A) baseline dispersive control with qubit rotations, cavity
displacements, and SNAP; (B) selective hybrid control with SQR-like primitives;
(C) ECD-like conditional displacement; (D) a minimal native chi-wait entangler
library; (E) waveform-level GRAPE; and (F) native SWAP-/sideband-style control
using Jaynes--Cummings and blue-sideband exchange primitives. During the study,
the local \texttt{cqed\_sim} copy was extended so that conditional displacement
and exchange-style native interactions are available as first-class synthesis
primitives rather than study-local wrappers.

The main result is that the best practical route is still a hybrid answer,
but the native-gate comparison is now substantially sharper. The exact logical
entangler is obtained most efficiently by native dispersive waiting plus fast
qubit rotations: \(F_{\mathrm{strict}}=0.9999999\) in \SI{256}{ns} with
negligible leakage. The best local cavity action still comes from the older
dispersive \(D\)-SNAP-\(D\) route, which reaches
\(F_{\mathrm{strict}}=0.9887\) with \(1.85\%\) average leakage. The new
SWAP-/sideband-style native family is a serious fast benchmark: it reaches
\(F_{\mathrm{strict}}=0.8784\) for the local task in \SI{440}{ns} and improves
on the selective SQR library, but it remains leakier and less exact than the
SNAP local gate and the native chi-wait entangler. GRAPE is useful as a pulse
benchmark, yet it does not beat the native entangler and only challenges the
local gate on speed, not fidelity. For this smallest Fock-encoded hybrid block,
the recommended route is therefore an augmented dispersive library
\(\{R_q, D, \mathrm{SNAP}, \text{native }\chi\text{-wait}\}\), with native
exchange primitives retained as a promising path for larger or less Fock-like
encodings.
\end{abstract}

\section{Introduction}
Circuit QED offers a natural hybrid architecture in which a nonlinear ancilla
qubit controls a high-coherence bosonic mode \cite{blais2004cqed,koch2007transmon}.
The resulting control question is no longer just whether a gate set is
universal in principle, but whether it is the best practical choice once gate
times, selectivity, leakage, calibration effort, and decoherence are all taken
seriously. That distinction is especially important for bosonic and hybrid
control because several paradigms compete on different operating principles.

The most mature dispersive route combines cavity displacements with
photon-number selective phase gates (SNAP) and ancilla rotations
\cite{heeres2015snap,krastanov2015universal,heeres2017logical}. Faster echoed
conditional displacement (ECD) strategies can outperform selective control in
weak-dispersive regimes \cite{eickbusch2022ecd}, and more recent
conditional-displacement variants such as CNOD push the same idea toward fast
structured multimode control \cite{diringer2024cnod}. A different native route
is to activate exchange or sideband interactions directly. Sideband control has
been part of the circuit-QED toolkit since the early qubit-photon entanglement
experiments of Wallraff \textit{et al.} \cite{wallraff2007sideband}, and more
recent work shows that one can rapidly transpose between dispersive and
resonant-like light-matter regimes \cite{valadares2024transposition} or use
sideband-enhanced SWAP control in weakly coupled multimode bosonic memories
\cite{huang2025fastsideband}. Finally, native activated nonlinear interactions
show what bosonic control can look like when the hardware itself supplies
stronger direct nonlinearities than a dispersive ancilla can offer
\cite{yang2024nativecubic}.

This study asks which of these ideas is best suited to the local simulator and
device assumptions already encoded in \texttt{cqed\_sim}. To keep the question
sharp, we start from the smallest hybrid logical space,
\(\{|g,0\rangle, |g,1\rangle, |e,0\rangle, |e,1\rangle\}\), and compare gate
libraries for two tasks: a local cavity basis change and a maximally entangling
hybrid unitary. Relative to the initial version of the study, the present
revision makes two important changes. First, native-gate families are compared
explicitly against selective ones. Second, the local \texttt{cqed\_sim} copy is
extended so ECD-like conditional displacement and SWAP-/sideband-style exchange
interactions can be benchmarked through the synthesis stack rather than only
discussed conceptually. The remainder of the paper presents the model and
benchmark, the candidate libraries, the comparative results, the validation
checks, and the resulting recommendation.

\section{System and Methods}

\subsection{Hamiltonian}
The static model is the dispersive transmon--cavity Hamiltonian supported by
\texttt{cqed\_sim},
\begin{equation}
\frac{H_0}{\hbar}
=
\omega_c a^\dagger a
+ \omega_q b^\dagger b
+ \frac{\alpha}{2} b^{\dagger 2} b^2
+ \chi\, a^\dagger a\, b^\dagger b
+ \frac{K}{2} a^{\dagger 2} a^2
+ H_{\chi'},
\end{equation}
where \(a\) is the cavity annihilation operator, \(b\) is the transmon lowering
operator, \(\alpha\) is the transmon anharmonicity, \(\chi\) is the leading
dispersive shift, \(K\) is the cavity self-Kerr term, and \(H_{\chi'}\) denotes
the first higher-order dispersive correction encoded in the model through
\texttt{chi\_higher=(\(\chi'\),)}. The direct-control benchmark uses the drive
Hamiltonian
\begin{equation}
\frac{H_d(t)}{\hbar}
=
\epsilon_c(t) a + \epsilon_c^\ast(t) a^\dagger
+ \epsilon_q(t) b + \epsilon_q^\ast(t) b^\dagger,
\end{equation}
with independently addressable cavity and qubit I/Q controls.

The native exchange benchmark is motivated by the effective red/blue-sideband
and Jaynes--Cummings interaction forms already supported by the simulator:
\begin{equation}
\frac{H_{\mathrm{JC}}}{\hbar}
=
g_{\mathrm{JC}}
\left(
a^\dagger \sigma_- e^{-i\phi_{\mathrm{JC}}}
+ a \sigma_+ e^{i\phi_{\mathrm{JC}}}
\right),
\end{equation}
\begin{equation}
\frac{H_{\mathrm{BS}}}{\hbar}
=
g_{\mathrm{BS}}
\left(
a^\dagger \sigma_+ e^{-i\phi_{\mathrm{BS}}}
+ a \sigma_- e^{i\phi_{\mathrm{BS}}}
\right),
\end{equation}
which represent ideal rotating-wave effective interactions rather than a full
calibrated pulse model. These are used at the synthesis layer as
\texttt{JaynesCummingsExchange} and \texttt{BlueSidebandExchange}.

The logical basis is the four-dimensional block
\begin{equation}
\mathcal{H}_{\mathrm{logical}}
=
\mathrm{span}\{
|g,0\rangle, |g,1\rangle, |e,0\rangle, |e,1\rangle
\}.
\end{equation}
Two benchmark targets are used:
\begin{equation}
U_H = I_q \otimes H_c,
\end{equation}
where \(H_c\) is the Hadamard on the cavity Fock pair \(\{|0\rangle,|1\rangle\}\),
and the cavity-controlled qubit flip
\begin{equation}
U_{\mathrm{CX}}
=
|0\rangle\langle 0|_c \otimes I_q
+ |1\rangle\langle 1|_c \otimes X_q.
\end{equation}

\subsection{Simulation Parameters}
\begin{table}[H]
  \centering
  \caption{Nominal system and control parameters used in the expanded benchmark.}
  \begin{tabular}{@{} l l l @{}}
    \toprule
    Parameter & Value & Unit \\
    \midrule
    Qubit frequency \(\omega_q/2\pi\) & 6.150 & GHz \\
    Cavity frequency \(\omega_c/2\pi\) & 5.241 & GHz \\
    Anharmonicity \(\alpha/2\pi\) & -255 & MHz \\
    Dispersive shift \(\chi/2\pi\) & -2.84 & MHz \\
    Higher-order dispersive shift \(\chi'/2\pi\) & -21 & kHz \\
    Cavity self-Kerr \(K/2\pi\) & -28 & kHz \\
    Cavity truncation \(n_{\mathrm{cav}}\) & 8 & Fock states \\
    Transmon truncation \(n_{\mathrm{tr}}\) & 2 & levels \\
    Qubit rotation duration & 40 & ns \\
    Displacement duration & 80 & ns \\
    Selective gate duration & 1100 & ns \\
    Conditional-displacement prior & 160 & ns \\
    Jaynes--Cummings exchange prior & 120 & ns \\
    Blue-sideband exchange prior & 120 & ns \\
    Native dispersive wait \(\pi/|\chi|\) & 176 & ns \\
    GRAPE amplitude bound \(|u_j|/2\pi\) & 40 & MHz \\
    Qubit \(T_1\) & 30 & \(\mu\)s \\
    Qubit \(T_2\) & 20 & \(\mu\)s \\
    Cavity \(T_1\) & 250 & \(\mu\)s \\
    \bottomrule
  \end{tabular}
\end{table}

\subsection{Computational Approach}
The study uses three layers of \texttt{cqed\_sim} functionality \cite{cqedsim}.
First, the device model is instantiated with
\texttt{DispersiveTransmonCavityModel} and a rotating \texttt{FrameSpec}.
Second, the decomposition study uses
\texttt{Subspace}, \texttt{UnitarySynthesizer}, \texttt{QubitRotation},
\texttt{Displacement}, \texttt{SNAP}, \texttt{SQR},
\texttt{FreeEvolveCondPhase}, \texttt{ConditionalDisplacement},
\texttt{JaynesCummingsExchange}, and \texttt{BlueSidebandExchange}. Third, the
pulse benchmark uses \texttt{build\_control\_problem\_from\_model},
\texttt{ModelControlChannelSpec}, \texttt{PiecewiseConstantTimeGrid},
\texttt{UnitaryObjective}, and \texttt{GrapeSolver}.

During this revision, the local \texttt{cqed\_sim} copy was extended so that
\texttt{ConditionalDisplacement}, \texttt{JaynesCummingsExchange}, and
\texttt{BlueSidebandExchange} are first-class sequence primitives in
\texttt{cqed\_sim.unitary\_synthesis}. They were also added to the fast-eval
path, exported through the package \texttt{\_\_init\_\_}, documented in the API
reference, and covered by dedicated unit tests. The remaining gap is not
operator availability but calibrated pulse export: the native primitives are
effective-interaction models, not yet waveform-bridge-backed pulse recipes.

Six libraries are compared:
\begin{align}
\mathcal{G}_A &= \{R_q, D, \mathrm{SNAP}\} \text{ with native }\chi\text{-wait allowed}, \\
\mathcal{G}_B &= \{R_q, D, \mathrm{SQR}\}, \\
\mathcal{G}_C &= \{R_q, D, \mathrm{CD}\}, \\
\mathcal{G}_D &= \{R_q, D, U_{\chi\text{-wait}}\}, \\
\mathcal{G}_E &= \{\text{piecewise qubit and cavity I/Q drives}\}, \\
\mathcal{G}_F &= \{R_q, U_{\mathrm{JC}}, U_{\mathrm{BS}}\}.
\end{align}
Here \(\mathrm{CD}\) is the first-class conditional-displacement primitive, and
\(U_{\mathrm{JC}}\) and \(U_{\mathrm{BS}}\) are the Jaynes--Cummings and
blue-sideband exchange primitives used to represent SWAP-/sideband-style native
control. Library \(D\) isolates the native dispersive entangler already present
in the hardware model.

For every candidate we report strict logical fidelity,
\begin{equation}
F_{\mathrm{strict}}
= \frac{1}{d^2} \left|\mathrm{Tr}\left(U_{\mathrm{target}}^\dagger U_{\mathrm{logical}}\right)\right|^2,
\end{equation}
with \(d=4\), the average leakage outside the logical block, total duration, and
a composite score
\begin{equation}
J
=
F_{\mathrm{logical}}
- 0.20\,L
- 0.02\,\frac{T}{\SI{1000}{ns}}.
\end{equation}
We also report a block-gauge fidelity, useful when a candidate reproduces the
correct entangling structure up to a removable cavity-block phase.

\section{Results}

\subsection{Goal 1: Local cavity action on the 2x2 block}
Figure~\ref{fig:local} and Table~\ref{tab:local} summarize the local cavity
benchmark. The best exact local action is still not the fastest pulse, but the
older dispersive \(D\)-SNAP-\(D\) route. Gate Set A reaches
\(F_{\mathrm{strict}}=0.9887\) with \(L=0.0185\) in \SI{1260}{ns}. This remains
the best strict and block-gauge score in the study and stays extremely stable
under perturbation replay and truncation checks.

The GRAPE benchmark is competitive but does not win the fidelity-first metric on
this target. The best nominal GRAPE local case is \(E_{\mathrm{local},320}\),
which reaches \(F_{\mathrm{strict}}=0.9618\) in \SI{320}{ns} with
\(L=0.0524\). This is a meaningful speed advantage, but it does not close the
accuracy gap to the structured SNAP-based decomposition.

The most important new local result is that the native SWAP-/sideband-style
family is a real intermediate candidate rather than a thought experiment.
Gate Set F reaches \(F_{\mathrm{strict}}=0.8784\) with \(L=0.0788\) in
\SI{440}{ns}. It therefore beats the selective SQR-based local library
(Gate Set B, \(F_{\mathrm{strict}}=0.8615\), \SI{2440}{ns}) on both fidelity and
speed, but still falls short of the SNAP baseline and of the best GRAPE pulse.
The ECD-like conditional-displacement library and the minimal native local probe
from Gate Set D both perform poorly for the strict Fock encoding, indicating
that neither simple displacement geometry nor bare chi-waiting alone is a good
local cavity-control primitive for this logical block.

\begin{figure}[H]
  \centering
  \includegraphics[width=0.95\linewidth]{../figures/fig1_local_tradeoffs.pdf}
  \caption{Local cavity-action benchmark for the target \(I_q\otimes H_c\). Point
  position gives duration and strict logical fidelity. Marker size decreases
  with leakage. The short black tick above a point indicates the block-gauge
  fidelity when it differs from the strict fidelity.}
  \label{fig:local}
\end{figure}

\begin{table}[H]
  \centering
  \caption{Best local-control candidates for the cavity Hadamard target.}
  \label{tab:local}
  \begin{tabular}{@{} l l l l l @{}}
    \toprule
    Library & Duration (ns) & \(F_{\mathrm{strict}}\) & Leakage & Score \(J\) \\
    \midrule
    A: \(D\)-SNAP-\(D\) & 1260 & 0.9887 & 0.0185 & 0.9598 \\
    E: GRAPE (\SI{320}{ns}) & 320 & 0.9618 & 0.0524 & 0.9449 \\
    F: native SWAP/sideband & 440 & 0.8784 & 0.0788 & 0.8538 \\
    B: selective hybrid & 2440 & 0.8615 & 0.1350 & 0.7857 \\
    D: native local probe & 336 & 0.5542 & 0.4747 & 0.4525 \\
    C: ECD-like CD & 440 & 0.4996 & 0.4535 & 0.4001 \\
    \bottomrule
  \end{tabular}
\end{table}

\subsection{Goal 2: Hybrid entangling unitary synthesis}
The entangler benchmark produces an even clearer answer. The native dispersive
conditional phase combined with two fast qubit rotations implements the target
\(\mathrm{CX}(c\rightarrow q)\) essentially exactly. Gate Set D reaches
\(F_{\mathrm{strict}}=0.9999999\) and \(L\approx 0\) in \SI{256}{ns}, and the
same sequence is recovered inside Gate Set A because the actual resource is the
native chi-wait primitive rather than SNAP itself. This is the dominant result
of the study: for the 2x2 Fock benchmark, the most important hybrid primitive is
already present in the native Hamiltonian and does not require either a
selective hybrid construction or a waveform-optimized pulse to become useful.

GRAPE again provides a useful reference but not a better answer. The best
nominal GRAPE entangler is \(E_{\mathrm{ent},400}\), which reaches
\(F_{\mathrm{strict}}=0.9458\) in \SI{400}{ns} with \(L=0.0854\). That is a
respectable benchmark, but it is neither faster nor cleaner than the native
dispersive entangler.

The expanded native comparison clarifies what SWAP-/sideband-style exchange can
and cannot do in this smallest logical block. Gate Set F reaches
\(F_{\mathrm{strict}}=0.7510\) with \(L=0.1834\) in \SI{320}{ns}. This is
substantially faster and slightly more faithful than the selective SQR entangler
on the strict metric, but it still falls well short of the exact native
conditional-phase route. Gate Set B remains physically interesting because a
single SQR entangler reaches block-gauge fidelity \(1.0\) but strict fidelity
\(1/\sqrt{2}\), which means the entangling structure is exactly right up to a
cavity-block phase that the current minimal selective library does not correct.
By contrast, the ECD-like library again performs poorly in the chosen encoding.

\begin{figure}[H]
  \centering
  \includegraphics[width=0.95\linewidth]{../figures/fig2_entangler_tradeoffs.pdf}
  \caption{Entangler benchmark for \(\mathrm{CX}(c\rightarrow q)\). The native
  dispersive wait plus qubit rotations dominates the Pareto front. Gate Set B
  achieves perfect block-gauge fidelity, revealing a missing local phase
  correction rather than a lack of entangling structure.}
  \label{fig:entangler}
\end{figure}

\begin{table}[H]
  \centering
  \caption{Best entangler candidates for the hybrid \(\mathrm{CX}(c\rightarrow q)\) target.}
  \label{tab:entangler}
  \begin{tabular}{@{} l l l l l l @{}}
    \toprule
    Library & Duration (ns) & \(F_{\mathrm{strict}}\) & \(F_{\mathrm{block}}\) & Leakage & Score \(J\) \\
    \midrule
    D: native chi-wait + \(R_q\) & 256 & 1.0000 & 1.0000 & 0.0000 & 0.9949 \\
    A: same entangler inside baseline library & 256 & 1.0000 & 1.0000 & 0.0000 & 0.9949 \\
    E: GRAPE (\SI{400}{ns}) & 400 & 0.9458 & 0.9467 & 0.0854 & 0.9207 \\
    F: native SWAP/sideband & 320 & 0.7510 & 0.7510 & 0.1834 & 0.7079 \\
    B: single SQR & 1100 & 0.7071 & 1.0000 & 0.0000 & 0.6851 \\
    C: ECD-like CD & 440 & 0.5000 & 0.5077 & 0.0000 & 0.4912 \\
    \bottomrule
  \end{tabular}
\end{table}

\subsection{Goal 3: Native vs Selective Control Routes}
The dedicated native-vs-selective comparison is shown in
Figure~\ref{fig:routes}. For the local cavity task, the best selective route is
still Gate Set A, the best native route is Gate Set F, and the best GRAPE pulse
lands between them: GRAPE is faster, the native SWAP/sideband route is more
interpretable, and SNAP remains the accuracy leader. For the entangler, the
picture changes completely. The decisive resource is native control: once the
chi-wait entangler is admitted explicitly, native control dominates both the
selective SQR library and the GRAPE reference on fidelity, leakage, and
duration. In that sense, the enlarged study sharpens the original conclusion.
The practical winner is not ``selective'' or ``native'' in the abstract, but a
split answer: selective dispersive control for the local cavity gate, native
interactions for the entangler.

The new Gate Set F also adds an important nuance. Native exchange interactions
are not merely future work; they already outperform the slower selective SQR
library on the present benchmark. What stops them from winning is not lack of
hybridity or speed, but logical exactness and leakage in the strict
\(\{|0\rangle,|1\rangle\}\) encoding. This is exactly the sort of regime where a
later cat/binomial study could plausibly reverse the ranking.

\begin{figure}[H]
  \centering
  \includegraphics[width=0.92\linewidth]{../figures/fig5_route_groups.pdf}
  \caption{Best selective, native, and GRAPE routes for each target. Selective
  control wins the local cavity gate through \(D\)-SNAP-\(D\), whereas native
  control wins the entangler through the chi-wait primitive. The SWAP-/sideband
  family is the strongest non-chi native local competitor but remains leakier
  than the selective baseline.}
  \label{fig:routes}
\end{figure}

\section{Validation}

\subsection{Sanity Checks}
The expanded study outputs satisfy several direct sanity checks. First, the
native dispersive entangler reproduces the target \(\mathrm{CX}(c\rightarrow q)\)
to machine precision on the logical block, both in the isolated native library
\(D\) and in the baseline library \(A\) when that same primitive is used.
Second, the SQR entangler reaches block-gauge fidelity \(1.0\), confirming that
the poor strict fidelity originates in a residual logical-phase mismatch rather
than in missing entangling structure. Third, the new SWAP-/sideband-style
library beats the selective SQR library on the local task and on strict
entangler fidelity, confirming that it is a meaningful native benchmark rather
than a contrived weak ansatz. Finally, the ECD-like library consistently
underperforms in the strict Fock encoding, which agrees with the observed large
leakage for the local cavity benchmark.

\subsection{Convergence Analysis}
The winning decomposition results are stable against cavity truncation. The
local \(D\)-SNAP-\(D\) sequence was rerun at \(n_{\mathrm{cav}}=6,8,10\), and
the strict fidelity stayed within \(1\times 10^{-7}\) of \(0.988734\). The
native entangler is even cleaner: its strict fidelity remains
\(0.9999999263\) at all three truncations with zero measured logical leakage.
Higher-budget reruns of both winning decomposition ansaetze returned exactly the
same nominal fidelities, so the decomposition conclusions are not limited by
optimizer noise.

The GRAPE reference required a more careful convergence interpretation. Coarse
single-seed grids were not reliable enough, so validation was performed in a
best-of-seed fine-grid regime at fixed duration. For the local benchmark at
\SI{480}{ns}, the best strict fidelities over three seeds were \(0.9713\),
\(0.9638\), and \(0.9843\) for 20, 24, and 32 time slices. For the entangler at
\SI{400}{ns}, the corresponding values were \(0.9737\), \(0.9397\), and
\(0.9317\) for 20, 24, and 28 slices. These variations are larger than for the
structured decompositions, but they are small enough to support the qualitative
benchmark conclusion that GRAPE does not beat the native entangler and only
challenges the SNAP route on local-gate duration.

\begin{figure}[H]
  \centering
  \includegraphics[width=0.95\linewidth]{../figures/fig3_robustness.pdf}
  \caption{Worst-case strict fidelity under \(\pm 2\%\) amplitude scaling,
  \(\pm 1\%\) dispersive-shift scaling, and \(+2\%\) duration stretch for the
  decomposition libraries. The winning Gate Set A local gate and the native
  chi-wait entangler remain the most stable candidates.}
  \label{fig:robustness}
\end{figure}

\subsection{Literature Comparison}
The observed hierarchy is consistent with the literature. The SNAP-based route
is indeed accurate and universal but slow because it depends on selective
operations \cite{heeres2015snap,krastanov2015universal}. The ECD-like route is
not disproven by this study; rather, the present results suggest that its main
advantages are likely to emerge in encodings more naturally aligned with
displacement geometry, as emphasized by the weak-dispersive ECD and CNOD
literature \cite{eickbusch2022ecd,diringer2024cnod}. The sideband/exchange
literature also matches the present native-gate outcome: these interactions are
naturally hybrid and can be very fast
\cite{wallraff2007sideband,valadares2024transposition,huang2025fastsideband},
but exact computational-basis matching is not automatic. That last point is an
inference from combining the literature with the present numerics: the papers
motivate exchange as a serious fast primitive, while the benchmark shows that
logical gauge control and leakage remain the main obstacles in the smallest
Fock block. Finally, the native-interaction perspective from
Ref.~\cite{yang2024nativecubic} matches the decisive entangler result: when the
hardware already supplies a useful nonlinear interaction, that native route can
dominate more elaborate selective or numerically optimized constructions.

\section{Discussion}
The central answer to the study question is that the best practical route is a
hybrid strategy that combines selective and native ideas rather than betting on
a single control family. For entanglement, the local device model already
provides the right primitive: wait under dispersive evolution for \(\pi/|\chi|\),
and dress that phase with fast qubit rotations. For local cavity control,
however, the older SNAP-plus-displacement construction still wins because it
gives exact, compact, and robust logical action in the chosen Fock block. Taken
together, those two facts point to an augmented dispersive library rather than
to a pure ECD route, a pure SWAP route, a pure SQR route, or a pure GRAPE route.

The enlarged native comparison also clarifies what the fast alternatives are
actually good for. The SWAP-/sideband-style library is not the winner here, but
it does beat the slower selective SQR library on both local-control score and
strict entangler fidelity. That matters because it shows the native-exchange
path is already competitive in the simulator, not just attractive in principle.
What holds it back in this particular benchmark is leakage and basis matching,
not lack of hybrid expressivity. By contrast, the ECD-like library remains
poorly matched to the strict \(\{|0\rangle,|1\rangle\}\) test even after the
study-local wrapper was replaced by a first-class synthesis primitive.

GRAPE remains valuable as a benchmark and as a possible source of compressed
pulse motifs, but its calibration burden and residual optimization sensitivity
make it less attractive as the primary answer when the native dispersive
entangler is already nearly perfect. The most important caveat is scaling. The
present ranking is specific to the smallest hybrid logical block. As the cavity
logical dimension grows, the cost of SNAP-based local control is expected to
rise, and this is exactly the regime in which ECD-, CNOD-, or exchange-inspired
control may become superior. The current recommendation should therefore be read
as the right answer for the 2x2 Fock benchmark, not as a universal answer for
all bosonic encodings.

\section{Conclusion}
The study goals were met. The local \texttt{cqed\_sim} codebase was audited for
hybrid-control capability, a literature-backed design memo was produced, the
synthesis layer was extended with first-class native primitives for conditional
displacement and SWAP-/sideband-style exchange, a reusable 2x2 logical
benchmark was implemented, multiple gate libraries were compared against common
targets, a GRAPE reference study was run, and the results were validated
against truncation and solver-budget checks.

For the present device model and logical encoding, the best practical route to
hybrid universal control is an augmented dispersive library:
\begin{equation}
\{R_q, D, \mathrm{SNAP}, \text{native }\chi\text{-wait}\}.
\end{equation}
The native dispersive conditional phase gives the best entangler by a wide
margin, while \(D\)-SNAP-\(D\) remains the best local cavity basis-change tool in
the strict 2x2 Fock benchmark. Native SWAP-/sideband-style control is now a
serious explicitly benchmarked candidate, and its performance suggests that it
should be revisited first when the study expands to larger or less Fock-like
encodings. GRAPE remains useful as a pulse benchmark and a possible compression
tool, but not as the preferred primary route in this first enlarged study.

\section{Limitations and Future Work}

\subsection{Known Limitations}
The first limitation is encoding scope. All conclusions were drawn for the
Fock-encoded cavity pair \(\{|0\rangle,|1\rangle\}\). This favors exact
low-photon-number phase control and may underestimate the value of
displacement-based fast control or exchange-style native control for cat-like or
other non-Fock logical bases.

The second limitation is the native implementation path. Although the study now
uses first-class \texttt{ConditionalDisplacement},
\texttt{JaynesCummingsExchange}, and \texttt{BlueSidebandExchange} primitives,
these are still ideal effective-interaction models in the synthesis layer. They
are not yet backed by a waveform bridge or a calibrated pulse-export workflow,
so the native results should be read as simulator-level gate-set benchmarks
rather than as fully compiled experimental schedules.

The third limitation is solver maturity on the GRAPE side. Although the GRAPE
benchmark was validated in a fine-grid best-of-seed regime, it still showed more
optimization sensitivity than the structured decompositions. The GRAPE numbers
therefore serve as a useful lower-bound benchmark, but not yet as a fully
converged optimum.

The fourth limitation is ancilla modeling. The decomposition study used a
two-level ancilla for tractability, which omits direct \(|f\rangle\)-manifold
leakage in the synthesis loop. That omission does not affect the main logical
comparison qualitatively, but it does limit how literally the present fidelities
should be interpreted experimentally, especially for sideband/exchange families.

\subsection{Suggested Improvements}
\begin{itemize}
  \item \textbf{[P1 | HIGH]} Add waveform-bridge / pulse-export support for
  \texttt{ConditionalDisplacement}, \texttt{JaynesCummingsExchange}, and
  \texttt{BlueSidebandExchange}. This is the main remaining framework gap
  between the present synthesis benchmark and full simulator-level replay.
  \item \textbf{[P1 | MEDIUM]} Repeat the entire benchmark on a cat-like
  two-dimensional cavity encoding. The present Fock benchmark likely
  disadvantages displacement-native and exchange-native control families.
  \item \textbf{[P2 | MEDIUM]} Upgrade the winning libraries to
  \(n_{\mathrm{tr}}=3\) noisy replay so direct \(|f\rangle\)-level leakage can be
  quantified.
  \item \textbf{[P2 | MEDIUM]} Add a local cavity block-phase correction
  primitive to complement SQR. The present single-SQR entangler achieved perfect
  block-gauge fidelity and may become an exact high-quality selective library
  once the phase mismatch is corrected cheaply.
  \item \textbf{[P2 | MEDIUM]} Add warm-started and uncertainty-aware GRAPE to
  reduce seed sensitivity and make the pulse benchmark more comparable to the
  structured native libraries.
  \item \textbf{[P3 | MEDIUM]} Extend the target set beyond the Hadamard and a
  single \(\mathrm{CX}\) to include random logical \(\mathrm{SU}(4)\) samples or
  a small hybrid Clifford benchmark.
\end{itemize}

\subsection{Open Questions}
The most important open question is whether the present winner remains dominant
once the cavity logical dimension grows. The study strongly suggests that the
local cavity-control bottleneck, not the entangler, is where the ranking could
change. A second open question is whether the SWAP-/sideband-style native family
would surpass the SNAP local gate in cat-like or binomial encodings, where
exchange structure may be more naturally aligned with the logical basis. A third
open question is whether an ECD- or CNOD-inspired local cavity primitive can
replace SNAP while preserving the exact native entangler. Finally, it remains to
be tested whether the best GRAPE pulses can be compressed into a short
interpretable library that is easier to calibrate experimentally.

\bibliographystyle{unsrtnat}
\bibliography{references}

\end{document}
